\documentclass[10pt,a4paper]{amsart}
\usepackage[margin=19mm]{geometry}
\usepackage{amsmath,amssymb,mathtools}
\usepackage[scr=rsfs,cal=euler]{mathalfa}
\usepackage{xfrac}
\usepackage[unicode,hidelinks,bookmarks=false]{hyperref}
\newcommand{\ie}{i.e.\ }
\newcommand{\cf}{cf.\ }
\newcommand{\resp}{respectively\ }
\newcommand{\Subsection}{\subsection}
\providecommand{\nobreakdash}{\nobreak\hbox{-}}
\setlength{\emergencystretch}{2em}
\multlinegap0pt
\usepackage{bm}
\usepackage{bbm}
\usepackage{dsfont}
\usepackage{xspace}
\usepackage{cancel}
\usepackage{mathtools}
\usepackage{cleveref}
\usepackage{amsthm}
\makeatletter
\@ifundefined{theorem}{\newtheorem{theorem}{Theorem}}{}
\makeatother
\title[Invariants of toric double determinantal rings]{Invariants of toric double determinantal rings}
\author{Jennifer Biermann, Emanuela De Negri, Oleksandra Gasanova, Asl{\i} Musapa\c{s}ao\u{g}lu, Sudeshna Roy}
\date{Source: arXiv:2506.22730v1 (CC BY 4.0)}
\begin{document}
\maketitle

\noindent\emph{Source and licence.} Invariants of toric double determinantal rings. Version arXiv:2506.22730v1 (CC BY 4.0), distributed under the Creative Commons Attribution 4.0 International licence, as recorded in the arXiv licence metadata for this exact version and shown on the abstract page https://arxiv.org/abs/2506.22730v1. Full rights notice: licenses/arxiv-2506.22730-CC-BY-4.0.txt.\par\smallskip

\noindent\textit{Editorial scope.} Counted unit: the Theorem whose source label is \texttt{thm: toric presentation} in the article (source0.tex lines 457--460, source label \texttt{thm: toric presentation}), delivered together with the article's own complete original proof of it (source0.tex lines 461--496, one \texttt{proof} environment of 36 lines). No mathematics is written, repaired or re-derived: statement and proof are the article's own text line for line. Required context retained uncounted: thm: mingens (lines 434--444). Every definition, display and label of those lines is retained unchanged. Omitted from this package and not redistributed: 1--433, 445--456, 497--997. Nothing inside a retained excerpt is omitted or altered: every retained body is delivered exactly as the article's lines are written, so no omission receipt of any kind accompanies this package. Citations: the retained excerpts carry no citation command at all, so no bibliography entry of the article is transcribed and no reference is redistributed. Reference adaptation: none was needed and none was made; every reference inside the delivered unit resolves inside it, so no unresolved reference or citation is produced. Printed slips kept literal: no printed word, symbol, hypothesis or number of the counted statement or of its proof was altered. Layout and class: the article's own class is \texttt{amsart} with packages including amscd, amsmath, amsthm, amssymb, verbatim, enumerate, color, amsfonts, latexsym, euscript, tikz-cd, geometry, enumitem, inputenc; the wrapper is the lane's pinned \texttt{amsart} editorial wrapper at 10pt on A4, which restates the theorem-like environments the retained text needs and adds the article's own macro definitions that the retained text needs (none), with extra wrapper packages (bm, bbm, dsfont, xspace, cancel, mathtools, cleveref). The delivered numbering is the unit's own. Assets: no retained span contains a figure environment, a graphics-inclusion command, a PDF-page inclusion, a table or a TikZ picture; no image file is copied into this package, the package declares no asset dependency and no image file is redistributed with it. Licence: version arXiv:2506.22730v1 of this article is distributed under the Creative Commons Attribution 4.0 International licence, as recorded in the arXiv licence metadata for this exact version and shown on the abstract page https://arxiv.org/abs/2506.22730v1. A full rights notice is delivered at licenses/arxiv-2506.22730-CC-BY-4.0.txt. Characters: every retained excerpt is pure ASCII, so the delivered engine, XeLaTeX with its default Latin Modern text font, reports no missing character.
\begin{theorem}
\label{thm: mingens}
 Let 
 \begin{align*}
   \mathcal M&=\{x_{ij_2}^{k_1}x_{ij_1}^{k_2}-x_{ij_1}^{k_1}x_{ij_2}^{k_2}\mid 1\le i\le m, \ 1\le j_1<j_2\le n, \ 1\le k_1<k_2\le r\},\\[1ex]
   \mathcal N&=\{x_{i_2j}^{k_1}x_{i_1j}^{k_2}-x_{i_1j}^{k_1}x_{i_2j}^{k_2}\mid 1\le i_1< i_2\le m, \  1\le j\le n, \ 1\le k_1<k_2\le r\},\\[1ex]
   \mathcal R&=\{x_{i_1j_2}^{k}x_{i_2j_1}^{k}-x_{i_1j_1}^{k}x_{i_2j_2}^{k}\mid 1\le i_1< i_2\le m, \  1\le j_1<j_2\le n, \ 1\le k\le r\},\\[1ex] 
   \mathcal T&=\{x_{i_1j_2}^{k_1}x_{i_2j_1}^{k_2}-x_{i_1j_1}^{k_1}x_{i_2j_2}^{k_2}, \ x_{i_2j_1}^{k_1}x_{i_1j_2}^{k_2}-x_{i_1j_1}^{k_1}x_{i_2j_2}^{k_2},\ x_{i_2j_2}^{k_1}x_{i_1j_1}^{k_2}-x_{i_1j_1}^{k_1}x_{i_2j_2}^{k_2}\mid \\& 1\le i_1< i_2\le m, \ 1\le j_1<j_2\le n, \ 1\le k_1<k_2\le r\}.
    \end{align*}
Then $I_{A_{mnr}}$ is minimally generated by $\mathcal M\cup \mathcal N\cup\mathcal R\cup\mathcal T$. 
\end{theorem}

\begin{theorem}
  \label{thm: toric presentation}
  With all the notation as above, one has $I_{mn}^r=I_{A_{mnr}}$.
  \end{theorem}

  \begin{proof} 
   First of all note that $I_{mn}^r\subseteq I_{A_{mnr}}$. Indeed, it is not so hard to directly check that any minor belongs to $\ker \phi$, where $\phi: R\to K[A_{mnr}]$ is given by $\phi(x_{ij}^k)=x_iy_jz_k$. In order to prove the other inclusion, it is enough to show that all the minimal generators of $I_{A_{mnr}}$, whose description we obtained in \Cref{thm: mingens}, belong to $I_{mn}^r$.
   \begin{enumerate}
       \item Any binomial in $\mathcal M$ is of the form $x_{ij_1}^{k_2}x_{ij_2}^{k_1}-x_{ij_1}^{k_1}x_{ij_2}^{k_2}$ for some $j_1<j_2$, $k_1<k_2$, thus it is equal, up to a sign,  to  $\begin{vmatrix}
    x_{ij_1}^{k_1} & x_{ij_2}^{k_1} \\[1ex]
     x_{ij_1}^{k_2}   &x_{ij_2}^{k_2}
    \end{vmatrix}$, which is a minor in $V.$ 
\item Any binomial in $\mathcal N$ is of the form $x_{i_1j}^{k_2}x_{i_2j}^{k_1}-x_{i_1j}^{k_1}x_{i_2j}^{k_2}$ for some $i_1<i_2$, $k_1<k_2$, thus it is equal, up to a sign, to  $\begin{vmatrix}
    x_{i_1j}^{k_1} & x_{i_1j}^{k_2} \\[1ex]
     x_{i_2j}^{k_1}   &x_{i_2j}^{k_2}
    \end{vmatrix}$, which is a minor in $H$.
\item Any binomial in $\mathcal R$ is of the form $x_{i_1j_2}^{k}x_{i_2j_1}^{k}-x_{i_1j_1}^{k}x_{i_2j_2}^{k}$ for some $i_1<i_2$, $j_1<j_2$, thus it is equal, up to a sign, to $\begin{vmatrix}
    x_{i_1j_1}^{k} & x_{i_1j_2}^{k} \\[1ex]
     x_{i_2j_1}^{k} &x_{i_2j_2}^{k}
    \end{vmatrix}$, which is a minor in the matrix $X_k$.
    \item Any binomial in $\mathcal T$ is of one of the following three forms, where $i_1<i_2$, $j_1<j_2$, $k_1<k_2$:
    \begin{itemize}
        \item $x_{i_1j_2}^{k_1}x_{i_2j_1}^{k_2}-x_{i_1j_1}^{k_1}x_{i_2j_2}^{k_2}$: up to a sign, this is equal to $\begin{vmatrix}
    x_{i_1j_1}^{k_1} & x_{i_1j_2}^{k_1} \\[1ex]
     x_{i_2j_1}^{k_2}   &x_{i_2j_2}^{k_2}
    \end{vmatrix}$, which is a minor in $V$.
        \item $x_{i_2j_1}^{k_1}x_{i_1j_2}^{k_2}-x_{i_1j_1}^{k_1}x_{i_2j_2}^{k_2}$: up to a sign, this is equal to $\begin{vmatrix}
    x_{i_1j_1}^{k_1} & x_{i_1j_2}^{k_2} \\[1ex]
     x_{i_2j_1}^{k_1}   &x_{i_2j_2}^{k_2}
    \end{vmatrix}$,  which is a minor in $H$.
        \item $x_{i_2j_2}^{k_1}x_{i_1j_1}^{k_2}-x_{i_1j_1}^{k_1}x_{i_2j_2}^{k_2}$: up to a sign, this is equal to
 $\begin{vmatrix}
    x_{i_1j_1}^{k_1} & x_{i_1j_2}^{k_1} \\[1ex]
     x_{i_2j_1}^{k_2}   &x_{i_2j_2}^{k_2}
    \end{vmatrix}+\begin{vmatrix}
    x_{i_1j_2}^{k_1} & x_{i_1j_1}^{k_2} \\[1ex]
     x_{i_2j_2}^{k_1}   &x_{i_2j_1}^{k_2}
    \end{vmatrix}$, which is the sum of a minor in $V$ and a minor in $H$.
    \end{itemize}       
   \end{enumerate}
  \end{proof}

\end{document}
